\chapter{Sheaf Cohomology and Leray}
\label{ch:sheaves}

We work with sheaves of modules on a ringed space \((X,\mathcal O_X)\).
An open subspace carries the restricted structure sheaf. For abelian
sheaves one can use the constant sheaf of rings \(\underline{\mathbb Z}\).
The calculations retain the module structures over the section rings.
In particular a map of ringed spaces induces restriction of scalars
on global sections even though its pushforward is already a functor of
categories of sheaves of modules.

\section{Cohomology and acyclicity on open sets}

The stalk functor at \(x\) is exact and has a right adjoint sending an
\(\mathcal O_{X,x}\)-module to its skyscraper sheaf at \(x\), with
sections equal to that module on opens containing \(x\) and zero
otherwise. The adjunction follows by sending a morphism into the
skyscraper to its value on germs at \(x\). Choose an embedding of each
stalk of \(\mathcal F\) into an injective \(\mathcal O_{X,x}\)-module.
The induced map into the product of the skyscrapers is monic, since
its projection at each \(x\) induces the chosen injection on the
\(x\)-stalk. Every factor is injective by the adjunction with an exact
functor, and the product is injective by extending one component at
a time. Thus the category of sheaves of modules has enough injectives.

For an open inclusion \(j:V\hookrightarrow X\), extension by zero
\(j_!\) is left adjoint to restriction and is exact. Its stalks are the
original stalks on \(V\) and zero outside \(V\), which proves exactness.
In particular the restriction of an injective sheaf to \(V\) is
injective. The adjunction identifies
\(\operatorname{Hom}(j_!\mathcal O_V,\mathcal F)\) with \(\mathcal F(V)\).
If \(V\subseteq W\) are open, the evident map of extensions by zero
of the unit modules is monic, as is seen on stalks. An injective
sheaf therefore has surjective restrictions between arbitrary opens.
A sheaf with this property is called flasque.

\begin{theorem}[Sections, flasque models, and change to abelian sheaves]
\label{res:H01}
For an open \(U\subseteq X\) define
\[
  H^i(U,\mathcal F)=H^i\bigl(\Gamma(U,I^\bullet)\bigr)
\]
from an injective resolution of \(\mathcal F\) on \(X\). This is also
the \(i\)-th right derived functor of sections of \(\mathcal F|_U\)
on \(U\). Its degree zero is \(\mathcal F(U)\), and a short exact
sequence of sheaves gives the natural cohomology sequence in
\(\mathcal O_X(U)\)-modules. Flasque sheaves are acyclic for sections
over every open. Taking underlying abelian sheaves identifies the
underlying groups of these cohomology objects with ordinary
cohomology of abelian sheaves.
\end{theorem}
\begin{proof}
The adjunction for \(j:U\hookrightarrow X\) shows that the restricted
injective resolution is an injective resolution on \(U\).
The degree-zero and long exact sequence claims follow from left
exactness of sections and Theorem~\ref{res:R09}.

To prove acyclicity, in \(0\to\mathcal F\to\mathcal G\to\mathcal H\to0\)
with \(\mathcal F\) flasque, every section of \(\mathcal H\) over an open
\(U\) lifts over \(U\). Partially lift it over open subsets of \(U\).
A chain of partial lifts glues over the union. If a maximal partial
lift omits a point of \(U\), the epimorphism of sheaves supplies a lift
near that point. On the overlap, the difference lies in \(\mathcal F\);
extend this difference to the new open and subtract it from the new lift.
The two lifts then glue, contradicting maximality. This proves
surjectivity on sections over every open. If \(\mathcal G\) is also
flasque, a section of \(\mathcal H\) lifts over its entire open
and its lift extends to any larger open. Hence \(\mathcal H\) is flasque.
Applied to an injective resolution of a flasque \(\mathcal F\), this
proves by induction that all successive cokernels are flasque and that
the resolution is exact after taking sections in positive degrees.

Forgetting scalars is exact on stalks and does not change restrictions.
An injective resolution of \(\mathcal F\) as an \(\mathcal O_X\)-module
becomes a flasque resolution of the underlying abelian sheaf, so it
computes the cohomology there by Theorem~\ref{res:R08}. The comparison
is natural because it is the canonical comparison through acyclic
models. On abelian sheaves one can also express sections on \(U\) as
\(\operatorname{Hom}(j_!\underline{\mathbb Z},-)\), so
Theorem~\ref{res:R11} gives the comparison with Ext from this free
abelian sheaf. See \cite[Tags 0712, 01E1, 01E3, 09SY]{stacks-sheaf}.
\end{proof}

\section{\v{C}ech cochains and the comparison spectral sequence}

\begin{definition}[\v{C}ech cochains for an indexed cover]\label{res:H02}
Let \(\mathcal U=(U_i)_{i\in I}\) be an open cover of \(U\).
Use all ordered tuples, including repetitions, and put
\[
  \check C^p(\mathcal U,\mathcal F)=
    \prod_{(i_0,\ldots,i_p)\in I^{p+1}}
       \mathcal F(U_{i_0}\cap\cdots\cap U_{i_p}).
\]
Its differential on the component \((i_0,\ldots,i_{p+1})\) is
the sum over \(a\) of \((-1)^a\) times the restriction of the
component with \(i_a\) omitted. The same formula defines the
complex for a presheaf of modules.
\end{definition}

Deleting positions \(a<b\) in opposite orders gives the same restriction
with signs \((-1)^{a+b-1}\) and \((-1)^{a+b}\), so \(d^2=0\).
The degree-zero augmentation from sections on \(U\) lands in the
kernel of the first differential. For a sheaf it is an isomorphism
onto that kernel, by the gluing axiom. All products and restrictions
are in \(\mathcal O_X(U)\)-modules via restriction of scalars.
See \cite[Sections 9--11]{stacks-sheaf}.

We need the following acyclicity calculation for the full, unnormalized
cochains. Write \(j_J:U_{i_0}\cap\cdots\cap U_{i_p}\to U\) for the open
inclusion indexed by a tuple \(J\). On \(U\) consider the augmented
chain complex with \(p\)-th term
\[
     \bigoplus_{J\in I^{p+1}}(j_J)_!\mathcal O_{U_J}
\]
and augmentation to \(\mathcal O_U\). The face maps on each summand
are the extension-by-zero inclusions, with the alternating signs.
At a point \(x\) the complex is the free chain module on tuples
of indices containing \(x\), over \(\mathcal O_{U,x}\). Choose one
index containing \(x\). Prepending it (also in degree \(-1\))
is a contraction of the augmented complex: the first deletion is
the identity and the other deletions cancel the prepending of the
old boundary. Stalks commute with these direct sums, so this proves
exactness of the augmented chain of sheaves. Apply
\(\operatorname{Hom}_{\mathcal O_U}(-,I')\) for an injective \(I'\).
Exactness is preserved contravariantly, and the adjunction for each
\(j_J\) gives exactly the product cochains above. Thus the
augmented \v{C}ech cochain complex is exact for each injective \(I'\).

\begin{theorem}[\v{C}ech-to-cohomology spectral sequence]\label{res:H03}
For a cover \(\mathcal U\) of \(U\) and a sheaf of modules \(\mathcal F\),
there is a first-quadrant spectral sequence
\[
 E_2^{p,q}=
 \check H^p\bigl(\mathcal U,\;V\longmapsto H^q(V,\mathcal F)\bigr)
 \quad\Longrightarrow\quad H^{p+q}(U,\mathcal F).
\]
It is natural in \(\mathcal F\) and induces a finite filtration of each
abutment group. The coefficient in the \v{C}ech expression is the
presheaf of local cohomology groups.
\end{theorem}
\begin{proof}
Let \(\mathcal F\to I^\bullet\) be an injective resolution. Form
\(D^{p,q}=\check C^p(\mathcal U,I^q)\), with commuting differentials
and the totalization of Proposition~\ref{res:B01}. The horizontal
augmentation by \(\Gamma(U,I^q)\) is exact for every \(q\), by the
preceding computation and exact open restriction. Therefore
\(\Gamma(U,I^\bullet)\to\operatorname{Tot}D\) is a quasi-isomorphism
by the augmented comparison in Theorem~\ref{res:B04}.

In the column filtration, vertical cohomology is the product of the
groups \(H^q(U_J,\mathcal F)\). Here products of modules are exact,
and \(I^\bullet|_{U_J}\) is an injective resolution. The horizontal
induced differential is the alternating restriction of these groups.
Consequently the second page is the displayed \v{C}ech term by
Proposition~\ref{res:B03}. First-quadrant support gives convergence and a
finite filtration by Theorems~\ref{res:B04} and~\ref{res:S01}.
Maps of coefficients lift to maps of injective resolutions and any
two lifts are homotopic by Theorem~\ref{res:R04}. On each column
the induced homotopy gives equality on vertical cohomology. The
total homotopy, multiplied by \((-1)^p\) in column \(p\), preserves
the column filtration and gives equality of the abutment maps.
Page functoriality now proves the stated naturality. See
\cite[Tags 01EP, 01ES]{stacks-sheaf}.
\end{proof}

\begin{corollary}[\v{C}ech comparison on acyclic intersections]
\label{res:H04}
If \(H^q(U_{i_0}\cap\cdots\cap U_{i_p},\mathcal F)=0\) for every
finite tuple and every \(q>0\), the natural comparison
\[
      \check H^p(\mathcal U,\mathcal F)\longrightarrow H^p(U,\mathcal F)
\]
is an isomorphism for each \(p\).
\end{corollary}
\begin{proof}
In Theorem~\ref{res:H03} only the \(q=0\) row survives. Its
identification with the target comes from the edge map, since the
target filtration has one nonzero graded layer in each total degree.
Concretely the comparison is induced by
\(\check C^\bullet(\mathcal U,\mathcal F)\to\operatorname{Tot}D\),
using the augmentation in vertical degree zero, followed on cohomology
by the inverse of the horizontal augmentation isomorphism.
It is therefore the usual comparison and natural in coefficients.
See \cite[Tag 01ET]{stacks-sheaf}.
\end{proof}

\begin{theorem}[Mayer--Vietoris for two opens]\label{res:H05}
For \(U=V\cup W\) there is a natural long exact sequence
\[
 \cdots\longrightarrow H^n(U,\mathcal F)
 \longrightarrow H^n(V,\mathcal F)\oplus H^n(W,\mathcal F)
 \xrightarrow{\operatorname{res}_V-\operatorname{res}_W}
 H^n(V\cap W,\mathcal F)
 \longrightarrow H^{n+1}(U,\mathcal F)\longrightarrow\cdots .
\]
\end{theorem}
\begin{proof}
For each injective \(I^q\) the sequence of sections is short exact.
Gluing identifies the kernel of the difference map. Surjectivity
follows by extending a section on \(V\cap W\) to \(V\), since \(I^q\)
is flasque. These short exact sequences form a short exact sequence
of complexes. Its ordinary cohomology sequence, with connecting map
given by lifting and taking the differential, is the displayed
sequence by Theorem~\ref{res:H01}. This convention agrees with the
triangle of Theorem~\ref{res:D04}; using the opposite order of \(V,W\)
reverses the sign of the difference map together with its boundary
identification. See \cite[Tag 01EB]{stacks-sheaf}.
\end{proof}

\section{Higher direct images and composition}

\begin{theorem}[Pushforward, local cohomology, and global sections]
\label{res:H06}
Let \(f:X\to Y\) be a map of ringed spaces. The sheaf \(R^q f_*\mathcal F\)
is the sheafification of the presheaf on \(Y\) sending \(V\) to
\(H^q(f^{-1}V,\mathcal F)\), with its natural \(\mathcal O_Y(V)\)-action.
If \(I\) is injective on \(X\), \(f_*I\) is acyclic for sections over
every open of \(Y\). There is a natural isomorphism
\[
  \operatorname{Res}\circ R\Gamma(X,-)
     \ \simeq\ R\Gamma(Y,-)\circ Rf_*
   \quad\text{on }D^+(\mathcal O_X\text{-}\mathrm{Mod}).
\]
Here \(\operatorname{Res}\) restricts scalars along
\(\mathcal O_Y(Y)\to\mathcal O_X(X)\).
\end{theorem}
\begin{proof}
The pushforward is additive and left exact. Its bounded right derived
functor is constructed in Theorem~\ref{res:R06}. For an injective
resolution \(I^\bullet\), cohomology sheaves of \(f_*I^\bullet\)
are obtained by sheafifying the presheaf cohomology, since the
comparison is an isomorphism on stalks by exactness of filtered
colimits of modules. Sections of \(f_*I^\bullet\) on \(V\) are sections
of \(I^\bullet\) on \(f^{-1}V\), giving the asserted local description
by Theorem~\ref{res:H01}. Pushforward preserves flasqueness since
the relevant restriction maps are those on inverse-image opens.
Thus \(f_*I\) is acyclic on every open by Theorem~\ref{res:H01}.
Apply the right composition comparison of Theorem~\ref{res:R12}
to \(f_*\) and \(\Gamma(Y,-)\). Their underived composite is
\(\operatorname{Res}\circ\Gamma(X,-)\). Restriction of scalars is
exact, and termwise restriction of the complex of sections on
an injective model computes the right derived composite.
This gives the displayed isomorphism, with the indicated coefficient
ring. See \cite[Tags 0713, 01E4, 01EX, 01EZ]{stacks-sheaf}.
\end{proof}

\begin{theorem}[Leray spectral sequence and two degenerations]
\label{res:H07}
For \(\mathcal F\) an \(\mathcal O_X\)-module, there is a first-quadrant
spectral sequence of \(\mathcal O_Y(Y)\)-modules
\[
 H^p(Y,R^q f_*\mathcal F)\ \Longrightarrow\
     \operatorname{Res} H^{p+q}(X,\mathcal F).
\]
For \(K\in D^+(\mathcal O_X\text{-}\mathrm{Mod})\), use \(H^q(Rf_*K)\)
on the second page and hypercohomology at the target. Here \(p\ge0\),
and \(q\in\mathbb Z\) is bounded below for each \(K\). Each abutment has a
finite filtration. For a sheaf \(\mathcal F\) the edge comparison gives
\[
\begin{array}{ll}
 H^n(Y,f_*\mathcal F)\ \simeq\ \operatorname{Res}H^n(X,\mathcal F)
    &\text{if all }R^q f_*\mathcal F=0\text{ for }q>0,\\[2pt]
 \operatorname{Res}H^n(X,\mathcal F)\ \simeq\
    \Gamma(Y,R^n f_*\mathcal F)
    &\text{if all }H^p(Y,R^q f_*\mathcal F)=0\text{ for }p>0 .
\end{array}
\]
For another map of ringed spaces \(g:Y\to Z\) the relative version has
second page \(R^p g_*(H^q(Rf_*K))\) and target
\(H^{p+q}(R(gf)_*K)\), with the same indexing convention.
\end{theorem}
\begin{proof}
Take \(F=f_*\) and \(G=\Gamma(Y,-)\) in
Theorem~\ref{res:G04}. The required acyclicity and the identification
of the composite are those of Theorem~\ref{res:H06}.
Corollary~\ref{res:G05} gives the first-quadrant case for a sheaf.
The two vanishing hypotheses leave only the row \(q=0\), or only the
column \(p=0\), respectively. The single surviving layer in each
total degree identifies the corresponding edge map with an isomorphism.

For the relative version, \(f_*I\) is flasque and is acyclic for
sections over every open of \(Y\). By the local description in
Theorem~\ref{res:H06} applied to \(g\), this implies
\(R^p g_*(f_*I)=0\) for \(p>0\). Thus the same hypothesis of the
derived composition theorem holds for \(g_*\).
Use \(g_*f_*=(gf)_*\) and Theorem~\ref{res:G04} again. The spectral
sequences and their edge comparisons are natural by the corresponding
constructions of Chapter~\ref{ch:grothendieck}.
See \cite[Tags 01F2, 01F4, 01F6]{stacks-sheaf}.
\end{proof}
